\section{Gestión del proyecto y entregables complementarios}\label{sec:gestion}

\subsection{Cronograma}
\begin{itemize}
\item{Cronograma correspondiente a las actividades del primer corte académico.}
\begin{table}[H]\centering\scriptsize\caption{Cronograma general del proyecto.}
\begin{tabularx}{\textwidth}{L{3.2cm}*{8}{C{0.75cm}}L{2.2cm}}\toprule
\textbf{Actividad}&\textbf{S1}&\textbf{S2}&\textbf{S3}&\textbf{S4}&\textbf{S5}&\textbf{S6}&\textbf{S7}&\textbf{S8}&\textbf{Responsable}\\\midrule
Requisitos&X&X&&&&&&&\campo{Nombre}\\
Diseño conceptual&&X&X&X&&&&&\campo{Nombre}\\
Prueba de comunicación&&&X&X&&&&&\campo{Nombre}\\
\campo{Actividad}&&&&&X&X&&&\campo{Nombre}\\\bottomrule
\end{tabularx}\end{table}

\item{Cronograma correspondiente a las actividades del segundo corte académico.}
\begin{table}[H]\centering\scriptsize
\caption{Cronograma general del proyecto — Corte 2 (Semanas 1--6).}
\label{tab:cronograma-corte2}
\begin{tabularx}{\textwidth}{L{4.0cm}*{6}{C{0.75cm}}L{4.3cm}}\toprule
\textbf{Actividad} & \textbf{S1} & \textbf{S2} & \textbf{S3} & \textbf{S4} & \textbf{S5} & \textbf{S6} & \textbf{Responsable}\\\midrule
Selección definitiva de componentes & X & & & & & & Karen Cecilia Villaronga, Fabián Ruiz Mateus\\
Actualización del diseño de ingeniería (P\&ID, CAD) & X & X & & & & & Yesid Gómez Rodríguez, Camilo Díaz Gamboa\\
Sistema eléctrico, alimentación y protecciones & & X & X & & & & Ferney Patiño Ojeda, Julián Santiago Jiménez, Karen Johanna García\\
Construcción e integración física & & X & X & X & & & Camilo Díaz Gamboa, Yesid Gómez Rodríguez, Ferney Patiño Ojeda\\
Modelado y control del proceso térmico & & & X & X & & & Fabián Ruiz Mateus, Yulieth Gutierrez Quintero\\
Programación de la secuencia CIP (GEMMA / Ladder) & & & X & X & X & & Julián Santiago Jiménez, Karen Cecilia Villaronga\\
Arquitectura centralizada y supervisión (PLC--HMI) & & & & X & X & & Camilo Díaz Gamboa, Yesid Gómez Rodríguez, Julian Santiago Jiménez\\
Comunicaciones industriales (verificación) & & & & & X & & Ferney Patiño Ojeda,  Camilo Díaz Gamboa\\
Integración y validación preliminar & & & & & X & X & Joseph Suárez Buitrago, Fabián Ruiz Mateus, Yulieth Gutierrez, Julian Santiago Jiménez\\
Documentación del avance & & & & & & X & Karen Johanna García, Karen Cecilia Villaronga\\\bottomrule
Repositorio organizado & X & X & X & X & X & X & Camilo Díaz Gamboa\\\bottomrule
\end{tabularx}\end{table}

\item{Cronograma correspondiente a las actividades del tercer corte académico.}
\begin{table}[H]\centering\scriptsize\caption{Cronograma general del proyecto.}
\begin{tabularx}{\textwidth}{L{3.2cm}*{8}{C{0.75cm}}L{2.2cm}}\toprule
\textbf{Actividad}&\textbf{S1}&\textbf{S2}&\textbf{S3}&\textbf{S4}&\textbf{S5}&\textbf{S6}&\textbf{S7}&\textbf{S8}&\textbf{Responsable}\\\midrule
Requisitos&X&X&&&&&&&\campo{Nombre}\\
Diseño conceptual&&X&X&X&&&&&\campo{Nombre}\\
Prueba de comunicación&&&X&X&&&&&\campo{Nombre}\\
\campo{Actividad}&&&&&X&X&&&\campo{Nombre}\\\bottomrule
\end{tabularx}\end{table}

\end{itemize}

\subsection{Presupuesto}
La Tabla~\ref{tab:presupuesto_ejecutado} presenta el presupuesto ejecutado en el segundo corte, a partir de las compras registradas para el prototipo, y la Tabla~\ref{tab:presupuesto_complementario}, los elementos prestados por la universidad, los donados y los instalados sin compra registrada.

% Gasto ejecutado del prototipo, a partir del formulario "Compras Proyecto Integrador 2026-2"
% (registros del 14/09/2026 al 15/10/2026). Soporte: archivos_complementarios/Compras_Proyecto_Integrador_2026-2.pdf
% Nombres y tags según el glosario de DECISIONES.md, la Tabla tab:instrumentacion y la Tabla tab:pid_elementos.
\begingroup
\footnotesize
\begin{longtable}{@{}L{3.5cm}L{5.75cm}C{0.9cm}>{\raggedleft\arraybackslash}p{1.75cm}L{1.9cm}@{}}
\caption{Presupuesto ejecutado: compras registradas para el prototipo.}
\label{tab:presupuesto_ejecutado} \\
\toprule
\textbf{Elemento} & \textbf{Especificación} & \textbf{Cant.} & \textbf{Valor} & \textbf{Compra}\\
\midrule
\endfirsthead
\multicolumn{5}{c}{\tablename\ \thetable{} -- continuación} \\
\toprule
\textbf{Elemento} & \textbf{Especificación} & \textbf{Cant.} & \textbf{Valor} & \textbf{Compra}\\
\midrule
\endhead
\midrule
\multicolumn{5}{r}{\textit{Continúa en la siguiente página}} \\
\endfoot
\bottomrule
\endlastfoot
\multicolumn{5}{@{}l}{\textbf{Calentamiento e instrumentación}} \\
\addlinespace[2pt]
E-101 Resistencia de inmersión & Fabricada a la medida, sin marca ni referencia comercial; 1500\,W teóricos, ha operado a $\approx$2000\,W con la tensión del laboratorio (2046\,W medidos a 126,3\,V). Se fija a la pared de TK-102 con contratuerca y empaques & 1 & \$120.000 & Julián Jiménez\\
LSL-101 / LSH-101 Flotador doble & Reed en acero inoxidable SUS304, 21,7\,cm, dos conmutadores (nivel bajo y alto) en TK-102 & 1 & \$55.500 & Julián Jiménez\\
Resistencia del lazo de TT-101 & 500\,$\Omega$; convierte la señal de 4--20\,mA del transmisor en 2--10\,V para la entrada analógica del PLC & 1 & \$350 & Fabián Ruiz\\
Resistencias del lazo de corriente y domicilio & Resistencias adicionales para el lazo de 4--20\,mA de TT-101, con envío & Lote & \$3.300 & Fabián Ruiz\\
\addlinespace[1pt]
\multicolumn{3}{@{}r}{\textit{Subtotal}} & \textit{\$179.150} & \\
\addlinespace
\multicolumn{5}{@{}l}{\textbf{Bombas y electroválvulas}} \\
\addlinespace[2pt]
P-102 Bomba de impulsión & Bomba periférica Pretul BOAP-1/2P3 de 1/2\,HP, 127\,V\,AC (placa), 60\,Hz, monofásica & 1 & \$114.776 & Julián Jiménez\\
P-103 Bomba de agitación sumergible & Modelo HJ-1141: 110--120\,V, 60\,Hz, 22\,W; caudal máximo 1000\,L/h; altura máxima 1,6\,m; inmersión máxima 2\,m; temperatura del líquido 5--35\,\textdegree C; protección IPX8. Permanece sumergida en TK-102 durante el calentamiento & 1 & \$70.000 & Camilo Díaz\\
P-101 Bomba dosificadora peristáltica & 12\,V\,DC, 80\,mA; caudal de 0--100\,mL/min; 0,1--100\,rpm en el cabezal (motor de 5000\,rpm); tubo de 2,5\,mm de diámetro interior y 4,7\,mm exterior; sentido de flujo según la polaridad; humedad relativa $<$80\,\%. Alimentada con el convertidor LM2596 & 1 & \$59.414 & Julián Jiménez\\
Electroválvulas & Electroválvulas de solenoide de 1/2\textquotedbl, normalmente cerradas, bobina de 24\,V\,DC; una unidad instalada como XV-102 (llenado); la segunda es servopilotada & 2 & \$42.990 & Julián Jiménez\\
\addlinespace[1pt]
\multicolumn{3}{@{}r}{\textit{Subtotal}} & \textit{\$287.180} & \\
\addlinespace
\multicolumn{5}{@{}l}{\textbf{Tubería y accesorios hidráulicos}} \\
\addlinespace[2pt]
Tubería de PVC & Tramos de 1\textquotedbl{} y 3/4\textquotedbl{} para el circuito y la tubería de prueba (REQ-05) & Lote & \$22.000 & Yesid Gómez\\
Accesorios de PVC & Codos de 1\textquotedbl{} y 3/4\textquotedbl, adaptadores macho de 1\textquotedbl, bujes y accesorios de 3/4\textquotedbl & Lote & \$20.200 & Yesid Gómez\\
Sellantes para tubería & Soldadura líquida para PVC, silicona y cinta de teflón & Lote & \$26.000 & Yesid Gómez\\
Abrazaderas & Abrazaderas de 1\textquotedbl{} y n.º 14 & Lote & \$4.300 & Yesid Gómez, Karen García\\
\addlinespace[1pt]
\multicolumn{3}{@{}r}{\textit{Subtotal}} & \textit{\$72.500} & \\
\addlinespace
\multicolumn{5}{@{}l}{\textbf{Tanques}} \\
\addlinespace[2pt]
TK-102 Tanque de mezcla & Tanque verde de plástico de 50\,L (20\,L por fase) & 1 & \$35.000 & Yesid Gómez\\
TK-103 Tanque de efluentes & Tanque azul metálico de 30\,L; previsto inicialmente como tanque de mezcla y reutilizado como tanque de efluentes & 1 & \$30.000 & Camilo Díaz\\
Montaje de E-101 en la pared de TK-102 & Contratuerca, empaques de sello, sellante y pincel para su aplicación & Lote & \$26.500 & Camilo Díaz\\
Acoples y perforación de los tanques & Acople de tanque, y sierra copa con su acople para las perforaciones & Lote & \$55.000 & Camilo Díaz\\
Reparación del tanque azul & Sellado de las perforaciones del tanque azul para su uso como TK-103 & Lote & \$15.200 & Ferney Patiño\\
\addlinespace[1pt]
\multicolumn{3}{@{}r}{\textit{Subtotal}} & \textit{\$161.700} & \\
\addlinespace
\multicolumn{5}{@{}l}{\textbf{Estructura}} \\
\addlinespace[2pt]
Estructura de los dos carros & Tubo cuadrado, perfil, láminas, ruedas y tornillería de la Tabla~\ref{tab:materiales_estructura} & Lote & \$500.000 & Ferney Patiño\\
Tornillos y bisagras & Tornillería y bisagras de montaje & Lote & \$6.000 & Camilo Díaz\\
Planchas de acrílico & Cerramiento del banco de relés del carro del PLC & Lote & \$12.000 & Yesid Gómez\\
\addlinespace[1pt]
\multicolumn{3}{@{}r}{\textit{Subtotal}} & \textit{\$518.000} & \\
\addlinespace
\multicolumn{5}{@{}l}{\textbf{Eléctrico y conexionado}} \\
\addlinespace[2pt]
SSR de la resistencia E-101 & SSR-60DA de 60\,A, entrada 3--32\,V\,DC, salida 24--380\,V\,AC, conmutación en cruce por cero, sin disipador adicional & 1 & \$41.500 & Julián Jiménez\\
Interruptor automático de E-101 & 15\,A, un polo, montaje en riel DIN & 1 & \$7.900 & Karen García\\
Convertidor reductor LM2596 & Entrada de 24\,V\,DC, salida ajustada a 12\,V\,DC para P-101 & 1 & \$6.000 & Karen García\\
Clavija y domicilio & Clavija de conexión a la red & 1 & \$8.000 & Fabián Ruiz\\
Conectores tipo banana & Para las extensiones de cableado hacia el banco del PLC & Lote & \$26.400 & Karen García\\
Terminales & Terminales de conexión (incluye terminales rojos) & Lote & \$9.000 & Karen García\\
Organización del cableado & Amarres y canaleta & Lote & \$11.500 & Karen García\\
\addlinespace[1pt]
\multicolumn{3}{@{}r}{\textit{Subtotal}} & \textit{\$110.300} & \\
\midrule
\multicolumn{3}{@{}r}{\textbf{Total ejecutado}} & \textbf{\$1.328.830} & \\
\end{longtable}
\endgroup


\begin{table}[H]
\centering
\footnotesize
\caption{Elementos del prototipo sin costo para el proyecto o pendientes de adquirir.}
\label{tab:presupuesto_complementario}
\begin{tabularx}{\textwidth}{@{}L{4.6cm}YL{2.2cm}>{\raggedleft\arraybackslash}p{1.5cm}@{}}
\toprule
\textbf{Elemento} & \textbf{Especificación} & \textbf{Origen} & \textbf{Valor}\\
\midrule
\multicolumn{4}{@{}l}{\textit{Prestados por la universidad}}\\
\addlinespace[2pt]
PLC & Siemens S7-1200 CPU 1214C AC/DC/RLY & Universidad & \$0\\
HMI & Siemens KTP600 Basic & Universidad & \$0\\
TE-101 / TT-101 Elemento y transmisor de temperatura & RTD PT100 de 3 hilos con transmisor SITRANS de 4--20\,mA (0--100\,\textdegree C) & Universidad & \$0\\
Switch industrial & Siemens CSM 1277 & Universidad & \$0\\
Raspberry Pi & Servidor local con Node-RED & Universidad & \$0\\
SSR de P-102 & 40\,A, sin disipador adicional & Universidad & \$0\\
Interruptor automático de las demás cargas & 7\,A (P-102, P-103, PLC y Raspberry Pi) & Universidad & \$0\\
Relé auxiliar de P-103 & Bobina de 24\,V\,DC & Universidad & \$0\\
Material del tablero & Cables, riel, bornera de 24\,V y 0\,V, y tomacorrientes & Universidad & \$0\\
Elementos de mando & Paro de emergencia tipo hongo, pulsadores, selector y pilotos & Universidad & \$0\\
Software & TIA Portal, licencia académica & Universidad & \$0\\
\addlinespace
\multicolumn{4}{@{}l}{\textit{Donados}}\\
\addlinespace[2pt]
Láminas de acrílico & Cerramiento de la zona de potencia del banco de relés & Donadas a un integrante & \$0\\
\addlinespace
\multicolumn{4}{@{}l}{\textit{Instalados sin compra registrada}}\\
\addlinespace[2pt]
TK-101 Tanque de detergente & Garrafa plástica de 1,8\,L, fijada en la parte inferior del carro de tanques & Sin compra registrada & \$0\\
\addlinespace
\multicolumn{4}{@{}l}{\textit{Pendientes de adquirir o de registrar}}\\
\addlinespace[2pt]
XV-103 Electroválvula de descarga & Solenoide de 24\,V\,DC, entre TK-102 y P-102; instalada en el tercer corte & \campo{Pendiente: origen} & Por registrar\\
HV-101 Válvula manual de vaciado & Vaciado de TK-103 al desagüe entre fases & Compra & Por cotizar\\
Manómetro con glicerina & 0--4\,bar, 1/4\textquotedbl{} (valor estimado en el primer corte) & Compra & \$25.000\\
\addlinespace
\multicolumn{4}{@{}l}{\textit{Contemplados en el primer corte y descartados en este corte (costo y diseño)}}\\
\addlinespace[2pt]
Válvulas de bola manuales & Regulación del lazo (1\textquotedbl) y desvío de tres vías & Descartadas & ---\\
Válvula antirretorno & Línea de dosificación & Descartada & ---\\
Termostato bimetálico & Corte de seguridad en serie con la resistencia; no se adquirió por presupuesto (se evaluará en el tercer corte) & Descartado & ---\\
\midrule
\multicolumn{3}{@{}r}{\textbf{Total estimado pendiente (sin XV-103 ni HV-101)}} & \textbf{\$25.000}\\
\bottomrule
\end{tabularx}
\end{table}


\paragraph{Comparación con el presupuesto preliminar.} El gasto ejecutado, de \$1.328.830, es \$497.786 menor que el presupuesto preliminar del primer corte, de \$1.826.616. La diferencia se debe principalmente a que varios elementos del diseño de referencia no se usaron en el prototipo, como el contactor, el guardamotor, el disipador y la ventilación del tablero, el motorreductor de agitación con su impulsor y su sello, el termostato bimetálico, las válvulas de bola, la válvula antirretorno y los módulos ESP32 y W5500 (Sección~\ref{sec:cambios_prototipo}). El termostato, las válvulas de bola y la válvula antirretorno se contemplaron en el primer corte como opciones posibles y se descartaron en este corte por costo y por simplificación del diseño. Aun así, el gasto ejecutado supera en \$328.830 el presupuesto objetivo de \$1.000.000. Las dos electroválvulas instaladas, XV-102 y XV-103, son las dos unidades iguales de esa compra, por lo que la de descarga no tuvo un costo adicional.


\subsection{Matriz de riesgos}
Para la evaluación de riesgos del proyecto se empleó una matriz de probabilidad e impacto, 
metodología semicuantitativa que permite priorizar los riesgos identificados en función de 
su criticidad. Cada riesgo fue calificado en dos dimensiones, utilizando una escala ordinal 
de 1 (muy bajo) a 5 (muy alto):

\begin{itemize}
    \item \textbf{Probabilidad}: posibilidad de que el riesgo se materialice durante el 
    desarrollo o la operación del sistema.
    \item \textbf{Impacto}: magnitud de las consecuencias sobre la seguridad del operador, 
    la integridad del equipo o el cumplimiento de los objetivos del proyecto, en caso de 
    que el riesgo se materialice.
\end{itemize}

El \textbf{nivel de riesgo} se obtiene como el producto de ambos factores 
(Nivel = Probabilidad $\times$ Impacto), lo que arroja valores entre 1 y 25. Estos valores 
se agrupan en tres categorías que orientan la urgencia del tratamiento:

\begin{itemize}
    \item \textbf{Riesgo bajo} (1--8): se acepta o se controla mediante medidas básicas 
    de rutina.
    \item \textbf{Riesgo medio} (9--12): requiere medidas de mitigación específicas, 
    documentadas e implementadas antes de la puesta en marcha.
    \item \textbf{Riesgo alto} (13--25): exige controles prioritarios y, de ser posible, 
    eliminación o sustitución de la fuente de peligro.
\end{itemize}

Los riesgos se agruparon en cinco categorías según su naturaleza: térmicos y de 
calentamiento, hidráulicos, de materiales y compatibilidad química, eléctricos y de 
control, y operativos/de proyecto. Para cada riesgo identificado se definió una medida de 
tratamiento orientada a reducir la probabilidad de ocurrencia, el impacto de sus 
consecuencias, o ambos. La Tabla~\ref{tab:matriz_riesgos} presenta la matriz preliminar 
resultante, donde se destacan en negrita los riesgos de mayor nivel (R-05 y R-21), 
correspondientes a la velocidad insuficiente del fluido de limpieza y al riesgo de 
sobrecosto presupuestal, respectivamente.
% Requiere en el preambulo: \usepackage{longtable}, \usepackage{booktabs}, \usepackage{array}, \usepackage{textcomp}
\begin{longtable}{@{}c p{4.5cm} ccc p{5.5cm}@{}}
\caption{Matriz preliminar de riesgos.}
\label{tab:matriz_riesgos} \\
\toprule
\textbf{ID} & \textbf{Riesgo} & \textbf{Prob.} & \textbf{Impacto} & \textbf{Nivel} & \textbf{Tratamiento} \\
\midrule
\endfirsthead

\multicolumn{6}{c}{\tablename\ \thetable{} -- continuación} \\
\toprule
\textbf{ID} & \textbf{Riesgo} & \textbf{Prob.} & \textbf{Impacto} & \textbf{Nivel} & \textbf{Tratamiento} \\
\midrule
\endhead

\midrule
\multicolumn{6}{r}{\textit{Continúa en la siguiente página}} \\
\endfoot

\bottomrule
\endlastfoot

\multicolumn{6}{@{}l}{\textit{Térmicos y de calentamiento}} \\
\addlinespace[2pt]
R-01 & Sobretemperatura del fluido de proceso & 2 & 5 & 10 & Corte de E-101 por programa del PLC al superar el umbral de seguridad, alarma alta por RTD y límite de consigna (el termostato bimetálico se descartó en este corte) \\
R-02 & Trabajo en seco de la resistencia & 3 & 4 & 12 & Enclavamiento por flotador de nivel mínimo como condición de habilitación \\
R-03 & Falla del SSR en cortocircuito & 2 & 5 & 10 & Pendiente para el tercer corte: protección independiente del PLC (el termostato bimetálico no se adquirió por presupuesto); mientras tanto, vigilancia de la temperatura durante el calentamiento \\
R-04 & Sobrecalentamiento del SSR por disipación insuficiente & 3 & 3 & 9 & Sin disipador adicional en este corte: vigilancia de la temperatura de los SSR durante el calentamiento; disipador por evaluar en el tercer corte \\
\addlinespace
\multicolumn{6}{@{}l}{\textit{Hidráulicos}} \\
\addlinespace[2pt]
R-05 & Velocidad inferior a 1,5\,m/s: limpieza no efectiva & 3 & 5 & \textbf{15} & Aforo volumétrico en puesta en marcha y manómetro con zona marcada \\
R-06 & Operación de la bomba contra válvula cerrada & 3 & 4 & 12 & Descarga de P-102 sin válvulas de cierre hacia TK-103 (la válvula de regulación del primer corte se descartó) \\
R-07 & Cavitación de la bomba & 3 & 3 & 9 & Succión inundada (TK-102 por encima de la succión de P-102) y línea corta \\
R-08 & Fuga hidráulica en acoples y uniones & 3 & 3 & 9 & Prueba hidrostática a 1,5\texttimes{} la presión de trabajo antes de cada campaña \\
R-09 & Formación de espuma: cavitación y arrastre & 4 & 3 & 12 & Detergente de bajo poder espumante; la bomba de agitación P-103 opera solo antes del calentamiento \\
R-10 & Obstrucción del filtro: caída de caudal & 4 & 2 & 8 & No aplica en este corte: el filtro Y no se usa en el prototipo \\
\addlinespace
\multicolumn{6}{@{}l}{\textit{Materiales y compatibilidad química}} \\
\addlinespace[2pt]
R-11 & Falla del sello mecánico por operar sobre 40\,\textdegree C & 4 & 3 & 12 & Consigna limitada a 35\,\textdegree C, por debajo del límite de 40\,\textdegree C de P-102, y enjuague final sin calentar \\
R-12 & Descincado del impulsor de latón & 3 & 3 & 9 & Detergente pH 8,0--10,5 sin amoniaco ni cáusticos clorados \\
R-13 & Degradación del tubo peristáltico: dosificación errática & 4 & 2 & 8 & Tubo Norprene, reemplazo cada 500\,h, verificación volumétrica mensual \\
R-14 & Residuo de detergente en la tubería tras el ciclo & 3 & 4 & 12 & Enjuague final de 20\,L en paso único; el método para verificar la ausencia de residuos se definirá en el tercer corte \\
\addlinespace
\multicolumn{6}{@{}l}{\textit{Eléctricos y de control}} \\
\addlinespace[2pt]
R-15 & Sobrecarga de los circuitos eléctricos & 3 & 3 & 9 & Circuitos separados: Q1 (15\,A) para la resistencia y Q2 (7\,A) para las demás cargas; los 16,2\,A medidos en la resistencia superan Q1 (pendiente para el tercer corte) \\
R-16 & Descarga eléctrica al operador & 2 & 5 & 10 & DDR 30\,mA, gabinete IP54, puesta a tierra y 24\,V SELV en campo \\
R-17 & Lectura de temperatura errónea & 3 & 3 & 9 & Verificación con termómetro patrón en dos puntos, referencia común a 2M \\
R-18 & Agotamiento de las salidas de relé por ciclado & 3 & 3 & 9 & Proporción de tiempo con periodo de 10\,s: la salida de relé del PLC conmuta el SSR en cada periodo, por lo que su desgaste se evaluará en el tercer corte \\
R-19 & Pérdida de red & 3 & 2 & 6 & Control local independiente en el PLC; Node-RED marca el sistema como fuera de línea \\
\addlinespace
\multicolumn{6}{@{}l}{\textit{Operativos y de proyecto}} \\
\addlinespace[2pt]
R-20 & Quemadura del operador por contacto con la resistencia o el fluido caliente & 3 & 4 & 12 & Aislamiento, etiquetado y enjuague frío enclavado antes del desacople \\
R-21 & Presupuesto excedido & 5 & 3 & \textbf{15} & Recortes priorizados y solicitud formal de ampliación \\
R-22 & Retraso por importación de componentes & 4 & 2 & 8 & Pedido con seis semanas de anticipación o sustitución local \\
\end{longtable}

\subsection{Repositorio digital}
\textbf{Enlace:} \url{\EnlaceRepositorio}

El repositorio \texttt{Proyecto\_CIP\_2026-2}, alojado en la cuenta \texttt{nicami13}, centraliza el código fuente en \LaTeX{} del informe, las figuras, los flujogramas, los programas del PLC y las actas del proyecto. A la fecha de este corte cuenta con más de 220 commits en la rama principal y con la participación de los 9 integrantes del equipo como colaboradores, lo que constituye evidencia directa de su participación en el control de versiones. El contenido se organiza por tipo de artefacto (secciones, tablas, figuras, flujogramas, referencias, anexos y actas) y no por integrante, para evitar conflictos de compilación y facilitar la trazabilidad de cada cambio; la administración general del repositorio recae principalmente en Karen Villaronga y Camilo Díaz Gamboa.

\subsection{Lista de verificación del segundo corte}
\begin{table}[H]\centering\small\caption{Lista de verificación del segundo corte.}
\label{tab:checklist_corte2}
\begin{tabularx}{\textwidth}{Y L{2.4cm} L{4.2cm}}\toprule
\textbf{Entregable}&\textbf{Estado}&\textbf{Ubicación/evidencia}\\\midrule
Secciones de planeación actualizadas al prototipo&Cumple&Secciones~\ref{sec:caracterizacion} a \ref{sec:plan_validacion}\\
Montaje físico: estructura e hidráulica&Cumple&Sección~\ref{sec:montaje}, p.\,\pageref{sec:montaje}\\
Instrumentación electrónica&Parcial: faltan P-101 y XV-103&Sección~\ref{sec:montaje_instrumentacion}, p.\,\pageref{sec:montaje_instrumentacion}\\
Programación&Parcial: versión preliminar&Sección~\ref{sec:montaje_programacion}, p.\,\pageref{sec:montaje_programacion}; Anexo~\ref{anx:programa_plc}\\
Comunicación&Parcial: evidencia prevista para el tercer corte&Sección~\ref{sec:montaje_comunicaciones}, p.\,\pageref{sec:montaje_comunicaciones}\\
Cambios entre el diseño y el prototipo&Parcial: motivos por confirmar&Sección~\ref{sec:cambios_prototipo}, p.\,\pageref{sec:cambios_prototipo}\\
Control: modelo experimental y controlador&Parcial: PI sin probar en lazo cerrado&Sección~\ref{sec:montaje_control}, p.\,\pageref{sec:montaje_control}\\
Pruebas&Parcial&Sección~\ref{sec:montaje_pruebas}, p.\,\pageref{sec:montaje_pruebas}\\
Evidencia en video&Pendiente&Sección~\ref{sec:montaje_video}, p.\,\pageref{sec:montaje_video}\\
Cronograma, presupuesto y riesgos&Cumple&Sección~\ref{sec:gestion}, p.\,\pageref{sec:gestion}\\
Repositorio organizado&Cumple&Sección~\ref{sec:gestion}\\
Declaración de IA&Cumple&Anexos\\
Reporte de similitud y artículo en inglés&Pendiente&Anexos\\\bottomrule
\end{tabularx}\end{table}


\subsection{Lista de verificación del tercer corte}
La Tabla~\ref{tab:checklist_corte3} resume el estado de las actividades del tercer corte (Tabla~\ref{tab:cronograma-corte3}) y de las pruebas de aceptación del proyecto.
\begin{table}[H]\centering\small\caption{Lista de verificación del tercer corte.}
\label{tab:checklist_corte3}
\begin{tabularx}{\textwidth}{Y L{3.4cm} L{3.6cm}}\toprule
\textbf{Actividad o prueba}&\textbf{Estado}&\textbf{Ubicación/evidencia}\\\midrule
Instalación de XV-103 y P-101&Cumple&Secciones~\ref{sec:montaje_hidraulico} y \ref{sec:montaje_instrumentacion}\\
Programa integrado: secuencia y control de temperatura&Parcial: falta la parada segura ante el paro de emergencia&Sección~\ref{sec:montaje_programacion}; Anexo~\ref{anx:programa_plc}\\
Tres ciclos completos&Cumple&VAL-01, Tabla~\ref{tab:resultados_pruebas}\\
Prueba continua de 30\,min&Pendiente&Sección~\ref{sec:montaje_pruebas}\\
Control en lazo cerrado y comparación con la simulación&Parcial: lazo cerrado cumple; falta la comparación&VAL-03, Sección~\ref{sec:montaje_control}\\
Tres consignas con perturbación controlada&Pendiente&Sección~\ref{sec:plan_validacion}\\
Prueba de tinta&Programada (26 de noviembre)&VAL-13\\
Supervisión remota (Node-RED, API e interfaz)&Cumple: periodo de 30\,s o menos&VAL-10, Sección~\ref{sec:montaje_comunicaciones}\\
Resiliencia ante pérdida de enlace&Pendiente&Sección~\ref{sec:montaje_comunicaciones}\\
Ciberseguridad de la supervisión&Pendiente: API sin HTTPS ni autenticación&Sección~\ref{sec:montaje_comunicaciones}\\
Alarmas detectadas y visualizadas&Cumple&VAL-09\\
Trazabilidad requisito--prueba--resultado&Parcial&Tabla~\ref{tab:resultados_pruebas}\\
Artículo final y reporte de similitud&Pendiente&Anexos\\
Video final de unos 5\,min&Pendiente&Sección~\ref{sec:montaje_video}\\
Presentación final y demostración&Pendiente&---\\
Repositorio organizado&Cumple&Sección~\ref{sec:gestion}\\\bottomrule
\end{tabularx}\end{table}

